\section{Model Training}
\label{sec:training}

\subsection{Training Setup}

The GestureTCN model is trained using PyTorch on a single GPU.
Table~\ref{tab:hyperparams} summarizes the key hyperparameters.

\begin{table}[ht]
\centering
\caption{Training hyperparameters.}
\label{tab:hyperparams}
\begin{tabular}{@{}ll@{}}
\toprule
\textbf{Hyperparameter} & \textbf{Value} \\
\midrule
Batch size & 32 \\
Maximum epochs & 300 \\
Optimizer & AdamW \\
Learning rate & $2 \times 10^{-3}$ \\
Weight decay & $1 \times 10^{-3}$ \\
LR scheduler & Cosine annealing ($\eta_{\min} = 10^{-5}$) \\
Gradient clipping & Max norm 2.0 \\
Early stopping patience & 40 epochs \\
\bottomrule
\end{tabular}
\end{table}

\subsection{Loss Function}

We use cross-entropy loss with label smoothing ($\alpha = 0.1$) and
inverse-frequency class weights:
\begin{equation}
  \mathcal{L} = -\sum_{c=1}^{C} w_c \, \tilde{y}_c \log p_c,
  \label{eq:loss}
\end{equation}
where $\tilde{y}_c = (1-\alpha) \cdot y_c + \alpha / C$ is the
smoothed label, $p_c$ is the predicted probability via softmax, and
$w_c$ is the class weight computed as $w_c = N / (C \cdot n_c)$ with
$n_c$ being the number of samples in class $c$. Label smoothing
prevents overconfident predictions and improves generalization, while
class weights address the moderate class imbalance in the dataset
(e.g., \texttt{swipe\_up} has the fewest training samples and
receives weight 1.259 versus 0.867 for \texttt{release}).

\subsection{Class Balancing}

In addition to weighted loss, we employ a \texttt{WeightedRandomSampler}
that assigns sampling probability inversely proportional to class
frequency. This ensures each class is sampled approximately equally
within each epoch, further mitigating class imbalance.

\subsection{Online Augmentation}

Beyond the offline augmentation described in
Section~\ref{sec:data}, the PyTorch \texttt{Dataset} applies
stochastic augmentation at each access:
\begin{itemize}[leftmargin=*]
  \item Coordinate jitter ($\sigma = 0.002$) with 50\% probability.
  \item Time warping with 30\% probability.
\end{itemize}
This means the model sees slightly different variants of each sample
every epoch, further improving generalization.

\subsection{Training Dynamics}

Training terminated after 68~epochs via early stopping, with the best
test accuracy of 89.66\% achieved at epoch~28.
Figure~\ref{fig:training_curves} shows the training and test loss/accuracy curves. The cosine annealing schedule
smoothly reduces the learning rate from $2\!\times\!10^{-3}$ to
$10^{-5}$, enabling fine-grained convergence in later epochs.

\begin{figure}[ht]
  \centering
  \includegraphics[width=0.95\linewidth]{training_curves}
  \caption{Training curves: loss (left) and accuracy (right) over
    68~epochs. Early stopping was triggered at epoch~68 with the best
    model saved at epoch~28.}
  \label{fig:training_curves}
\end{figure}
